% \begin{UPSTIinfor}{Rappel : le pseudo-code}
%     Comme au TD précédent, on décrira parfois l'idée d'un programme en \textbf{pseudo-code}, indépendamment du C : \verb|variable <- expression| pour une affectation, \verb|SI ... ALORS ... SINON ... FIN SI| pour un aiguillage, \verb|AFFICHER ...| pour un affichage. On y ajoute ici la notation pour une répétition :
%     \begin{verbatim}
% TANT QUE condition FAIRE
%     instruction(s)
% FIN TANT QUE
%     \end{verbatim}
%     qui se lit \og tant que la condition est vraie, on répète les instructions \fg, et qui est l'équivalent du bloc Scratch \og répéter jusqu'à... \fg{} (à la condition inversée près).
% \end{UPSTIinfor}

\begin{UPSTIinfor}{Trace d'une boucle \texttt{while}}
    Tracer l'exécution d'une boucle, c'est noter, tour par tour, la valeur de chaque variable et le résultat du test de la condition, jusqu'à ce que la boucle s'arrête.
    \begin{lstlisting}[language=C]
int i = 0;
while (i < 3)
{
    printf("Tour %d\n", i);
    i = i + 1;
}
\end{lstlisting}
    \begin{center}
        \begin{tabular}{|c|c|c|}
            \hline
            \textbf{Avant le test} & \texttt{i < 3} ? & \textbf{Affiché, puis \texttt{i} devient} \\
            \hline
            i = 0 & vrai & \og Tour 0\fg, i = 1 \\
            \hline
            i = 1 & vrai & \og Tour 1\fg, i = 2 \\
            \hline
            i = 2 & vrai & \og Tour 2\fg, i = 3 \\
            \hline
            i = 3 & faux & (rien, la boucle s'arrête) \\
            \hline
        \end{tabular}
    \end{center}
\end{UPSTIinfor}

\section{Revoir les bases des boucles.}
\begin{UPSTIexercice}{Trace d'une boucle \texttt{while}}
    On exécute le code suivant.
    \begin{lstlisting}[language=C]
int compteur = 5;
while (compteur > 0)
{
    printf("Reste %d\n", compteur);
    compteur = compteur - 2;
}
printf("Fin, compteur = %d\n", compteur);
\end{lstlisting}
    \UPSTIquestion{Compléter le tableau de trace ci-dessous, représentant ce qu'affiche le programme. Lorsque l'on se trouve à l'intérieur de la boucle, indiquer la valeur de \texttt{compteur} \textbf{avant} le test de la condition ainsi que le résultat de ce test.}

    \begin{center}
        \begin{tabular}{|c|c|p{4cm}|}
            \hline
            \textbf{Valeur de \texttt{compteur} avant le test} & \textbf{\texttt{compteur > 0} ?} & \textbf{Affiché} \\
            \hline
             & & \\[0.3cm]
            \hline
             & & \\[0.3cm]
            \hline
             & & \\[0.3cm]
            \hline
             & & \\[0.3cm]
            \hline
        \end{tabular}
    \end{center}
    \UPSTIquestion{Quelle est la dernière ligne affichée par le programme (après la sortie de la boucle) ?}
\end{UPSTIexercice}

\begin{UPSTIprofOnlyEnv}
    \begin{UPSTIcorrectionP}{Trace d'une boucle \texttt{while}}
        \begin{center}
            \begin{tabular}{|c|c|c|}
                \hline
                \textbf{\texttt{compteur} avant le test} & \textbf{\texttt{compteur > 0} ?} & \textbf{Affiché} \\
                \hline
                5 & vrai & \og Reste 5\fg, compteur devient 3 \\
                \hline
                3 & vrai & \og Reste 3\fg, compteur devient 1 \\
                \hline
                1 & vrai & \og Reste 1\fg, compteur devient -1 \\
                \hline
                -1 & faux & (rien, sortie de boucle) \\
                \hline
            \end{tabular}
        \end{center}
        Dernière ligne affichée : \texttt{Fin, compteur = -1}. Le compteur passe directement de 1 à -1 sans jamais valoir exactement 0 : la condition d'arrêt \texttt{compteur > 0} n'a pas besoin que \texttt{compteur} atteigne 0 pile pour devenir fausse.
    \end{UPSTIcorrectionP}
\end{UPSTIprofOnlyEnv}


\begin{UPSTIwarning}{La boucle infinie}
    Une boucle \texttt{while} dont la condition reste \textbf{toujours vraie} ne s'arrête jamais : c'est une \textbf{boucle infinie}. C'est presque toujours une erreur, le plus souvent parce que l'on a oublié de faire évoluer, à l'intérieur de la boucle, la variable testée par la condition.
\end{UPSTIwarning}

\begin{UPSTIexercice}{Trouver et corriger une boucle infinie}
    Le code suivant devrait afficher les nombres de 0 à 4, mais il ne s'arrête jamais lorsqu'on l'exécute.
    \begin{lstlisting}[language=C]
int i = 0;
while (i < 5)
{
    printf("%d\n", i);
}
\end{lstlisting}
    \UPSTIquestion{Expliquer précisément pourquoi ce code ne s'arrête jamais : que vaut \texttt{i} à chaque tour de boucle ?}
    \UPSTIquestion{Corriger le code pour qu'il affiche bien les nombres de 0 à 4 puis s'arrête.}
\end{UPSTIexercice}

\begin{UPSTIprofOnlyEnv}
    \begin{UPSTIcorrectionP}{Trouver et corriger une boucle infinie}
        La variable \texttt{i} vaut \texttt{0} à chaque tour de boucle : rien, dans le corps de la boucle, ne modifie \texttt{i}. La condition \texttt{i < 5} reste donc toujours vraie, et la boucle ne s'arrête jamais.
        \begin{lstlisting}[language=C]
int i = 0;
while (i < 5)
{
    printf("%d\n", i);
    i = i + 1;
}
\end{lstlisting}
    \end{UPSTIcorrectionP}
\end{UPSTIprofOnlyEnv}



\begin{UPSTIexercice}{Repérer les trois parties d'un \texttt{for}}
    On considère le code suivant.
    \begin{lstlisting}[language=C]
for (int compteurPieces = 1; compteurPieces <= 6; compteurPieces = compteurPieces + 1)
{
    printf("Piece numero %d fabriquee\n", compteurPieces);
}
\end{lstlisting}
    \UPSTIquestion{\textbf{Entourer} l'\textbf{initialisation}, \textbf{souligner} la \textbf{condition} et \textbf{encadrer} la \textbf{mise à jour} dans l'en-tête du \texttt{for} (les trois parties séparées par des \texttt{;}).}
    \UPSTIquestion{Combien de fois le bloc d'instructions est-il exécuté ? Justifier en indiquant la première et la dernière valeur prise par \texttt{compteurPieces}.}
    \UPSTIquestion{Réécrire ce même programme avec une boucle \texttt{while}, en faisant apparaître explicitement, comme dans le \texttt{for}, l'initialisation avant la boucle et la mise à jour dans le corps de la boucle.}
\end{UPSTIexercice}

\begin{UPSTIprofOnlyEnv}
    \begin{UPSTIcorrectionP}{Repérer les trois parties d'un \texttt{for}}
        Initialisation : \texttt{int compteurPieces = 1}. Condition : \texttt{compteurPieces <= 6}. Mise à jour : \texttt{compteurPieces = compteurPieces + 1}.

        Le bloc s'exécute \textbf{6 fois}, pour \texttt{compteurPieces} valant successivement 1, 2, 3, 4, 5, 6 (la boucle s'arrête dès que \texttt{compteurPieces} vaut 7, car la condition devient fausse).

        \begin{lstlisting}[language=C]
int compteurPieces = 1;
while (compteurPieces <= 6)
{
    printf("Piece numero %d fabriquee\n", compteurPieces);
    compteurPieces = compteurPieces + 1;
}
\end{lstlisting}
    \end{UPSTIcorrectionP}
\end{UPSTIprofOnlyEnv}


\begin{UPSTIexercice}{Prédire l'affichage d'un \texttt{for}}
    On exécute le programme suivant.
    \begin{lstlisting}[language=C]
for (int i = 10; i > 0; i = i - 3)
{
    printf("i = %d\n", i);
}
printf("Boucle terminee\n");
\end{lstlisting}
    \UPSTIquestion{Prédire, ligne par ligne, tout ce que ce programme affiche.}
\end{UPSTIexercice}

\begin{UPSTIprofOnlyEnv}
    \begin{UPSTIcorrectionP}{Prédire l'affichage d'un \texttt{for}}
        \begin{lstlisting}[language=bash,style=console]
i = 10
i = 7
i = 4
i = 1
Boucle terminee
\end{lstlisting}
    \end{UPSTIcorrectionP}
\end{UPSTIprofOnlyEnv}


\begin{UPSTIwarning}{Le piège du \og un cran \fg{} (off-by-one)}
    Utiliser \texttt{<} au lieu de \texttt{<=} (ou inversement) dans la condition d'une boucle est l'une des erreurs les plus fréquentes en programmation : elle ne provoque \textbf{aucune erreur de compilation}, mais fait exécuter la boucle une fois de trop, ou une fois de moins, que prévu.
\end{UPSTIwarning}

\begin{UPSTIexercice}{Un cran de trop, ou de trop peu}
    On veut afficher exactement les nombres de 1 à 5 (bornes incluses). Un étudiant propose le code suivant :
    \begin{lstlisting}[language=C]
for (int i = 1; i < 5; i = i + 1)
{
    printf("%d\n", i);
}
\end{lstlisting}
    \UPSTIquestion{Ce code affiche-t-il bien les nombres de 1 à 5 ? Sinon, préciser.}
    \UPSTIquestion{Corriger le code.}
\end{UPSTIexercice}

\begin{UPSTIprofOnlyEnv}
    \begin{UPSTIcorrectionP}{Un cran de trop, ou de trop peu}
        Ce code affiche seulement \texttt{1, 2, 3, 4} : le nombre \texttt{5} \textbf{manque}, car la condition \texttt{i < 5} devient fausse dès que \texttt{i} atteint 5, avant même l'exécution du bloc pour cette valeur.
        \begin{lstlisting}[language=C]
for (int i = 1; i <= 5; i = i + 1)
{
    printf("%d\n", i);
}
\end{lstlisting}
    \end{UPSTIcorrectionP}
\end{UPSTIprofOnlyEnv}


\begin{UPSTIexercice}{Traduire un \texttt{while} en \texttt{for}}
    On considère le code suivant, qui affiche une jauge de charge de batterie par paliers de 20\%.
    \begin{lstlisting}[language=C]
int charge = 0;
while (charge <= 100)
{
    printf("Charge : %d %%\n", charge);
    charge = charge + 20;
}
\end{lstlisting}
    \UPSTIquestion{Identifier, dans ce code, ce qui correspond à l'\textbf{initialisation}, la \textbf{condition} et à la \textbf{mise à jour} d'une boucle \texttt{for}.}
    \UPSTIquestion{Réécrire ce programme en utilisant une boucle \texttt{for} unique, produisant exactement le même affichage.}
\end{UPSTIexercice}

\begin{UPSTIprofOnlyEnv}
    \begin{UPSTIcorrectionP}{Traduire un \texttt{while} en \texttt{for}}
        Initialisation : \texttt{charge = 0} (avant la boucle). Condition : \texttt{charge <= 100}. Mise à jour : \texttt{charge = charge + 20} (à la fin du corps de la boucle).
        \begin{lstlisting}[language=C]
for (int charge = 0; charge <= 100; charge = charge + 20)
{
    printf("Charge : %d %%\n", charge);
}
\end{lstlisting}
    \end{UPSTIcorrectionP}
\end{UPSTIprofOnlyEnv}
